\section{Finite source controls and a two-state readout pointer}
\label{sec:chvisd-finite-readout}

This construction specifies a source-selective receiving model. Its
controls act on the original sixteen-state source and a finite pointer;
the complete continuum field remains active. It supplies a finite
mathematical implementation of the stated source effects with an explicit
drift and switching-work budget. The actuator implementing these source
operators, its timing and calibration, and any autonomous replacement for
the driven controls remain physical apparatus assignments.

\subsection{The retained source energy is enough for short pulses}
\label{sec:chvisd-readout-form-energy}

In the declared prepared lane retain
\begin{equation}
\begin{gathered}
H_0=\epsilon eI+\epsilon S+K,\\
 K=H_f+H_{\rm int}+cI\geq0,\\
 c=2g^2\mu.
\end{gathered}
\label{eq:chvisd-readout-source-generator}
\end{equation}
Here $S$ is the original Swap restriction, $\|S\|=1$, and the even energy,
both physical polarizations, both counterterms, hard dressing, soft vacuum,
and all field complements are retained. On the original common form
domain the annihilation form estimate gives
\begin{equation}
 |\langle H_{\rm int}\rangle|
      \leq2\sqrt{c\langle H_f\rangle}
      \leq\langle H_f\rangle+c.
\end{equation}
Consequently, with the lane scalar removed only from the comparison phase,
\begin{equation}
 \begin{aligned}
 L&=H_0-\epsilon eI+\epsilon I\geq0,\\
 L&\leq2\epsilon I+2H_f+2cI
                   \quad\hbox{as quadratic forms}.
 \end{aligned}
 \label{eq:chvisd-readout-positive-form}
\end{equation}
The scalar $\epsilon e-\epsilon$ is still part of the physical generator
and its elapsed phase; removing its common phase from this comparison
does not subtract its energy or any relative route phase.

Precisely, the positivity and scalar subtraction in
Eq.~\eqref{eq:chvisd-readout-positive-form} apply on the selected invariant
lane tensored with the full Fock space and pointer. They are not a global
all-sixteen-state positivity claim after subtracting one lane's $e$.
The actual $H_0$ and bounded source controls are defined on all sixteen
source states. Swap and $P_{ba}$ preserve each invariant source lane, so
the finite-energy proof restricts to the reference input lane without
discarding the other states from the physical dynamics.

Suppose the declared reference source--field input to the readout has
$\langle H_f\rangle\leq F$ at its declared entry age. Adjoin the blank
finite pointer. Let $V$ be an ideal source--pointer unitary and $B$ a
self-adjoint contraction, both preserving the selected lane and commuting
with $H_f$. The controls used below have these properties. Then
for $\xi=V\psi$ the form bounds imply
\begin{equation}
\begin{gathered}
\langle\xi,L\xi\rangle\leq C,\\
 \langle B\xi,L B\xi\rangle\leq C,\\
 C=2\epsilon+2F+2c.
\end{gathered}
\label{eq:chvisd-readout-common-C}
\end{equation}
The second expression is not normalized: $\|B\xi\|\leq1$ and
$\langle B\xi,H_fB\xi\rangle\leq F$. Commutation with $H_f$ preserves
its form domain and is used here; preservation of $D(H_0)$ by $B$ or $V$
is neither assumed nor needed. A mixed input is covered by a purification,
with the same expectation bounds and an untouched reference.

The retained source estimate can be inserted before selecting pulse times.
For example, if the declared input-age bound is
\begin{equation}
 F\leq2c+2K_{\rm prep}+2\epsilon E(T_0),
\end{equation}
then one may use the outward source-bound ceiling
\begin{equation}
 C\leq C^{\rm up}
   =2\epsilon+6c+4K_{\rm prep}+4\epsilon E(T_0).
 \label{eq:chvisd-readout-source-C-upper}
\end{equation}
This substitution inherits the preparation, input-age window and
definitions of its source bound. It does not assign a new measured
accuracy or an energy second moment.

\subsection{A bounded-control drift lemma}
\label{sec:chvisd-bounded-pulse-lemma}

For $L\geq0$ and any vector in its form domain, the spectral theorem gives
\begin{equation}
 \|(e^{-itL}-I)\xi\|^2
   =\int 2[1-\cos(tE)]\,d\mathcal P_\xi(E)
   \leq2|t|\langle\xi,L\xi\rangle.
 \label{eq:chvisd-free-form-increment}
\end{equation}
Thus a finite first energy form suffices, even when $L\xi$ is not a
Hilbert-space vector.

Let a pulse of duration $\tau>0$ have actual generator
\begin{equation}
\begin{gathered}
H(t)=L+\lambda(t)B,\\
 A_{\rm abs}=\int_0^\tau|\lambda(t)|\,dt<\infty,
\end{gathered}
\label{eq:chvisd-readout-driven-generator}
\end{equation}
with $B$ as above. The finite pointer is degenerate during this modelled
control, and $\lambda(t)$ is real. In the interaction picture put
\begin{equation}
\begin{gathered}
B_H(t)=e^{itL}Be^{-itL},\\
 V(t)=\exp\!\left[-iB\int_0^t\lambda(u)\,du\right].
\end{gathered}
\end{equation}
The ideal $V(t)$ commutes with $H_f$. For
$\xi_t=V(t)\psi$, split
\begin{equation}
 \begin{aligned}
 (B_H(t)-B)\xi_t
   ={}&e^{itL}B(e^{-itL}-I)\xi_t\\
      &+(e^{itL}-I)B\xi_t.
 \end{aligned}
\end{equation}
Equations~\eqref{eq:chvisd-readout-common-C} and
\eqref{eq:chvisd-free-form-increment} therefore give
\begin{equation}
\begin{gathered}
\|(B_H(t)-B)V(t)\psi\|
          \leq2\sqrt{2tC},\\
 0\leq t\leq\tau.
\end{gathered}
\label{eq:chvisd-interaction-control-drift}
\end{equation}
All free factors in this calculation are unitary on the Hilbert space.
No commutator $[L,B]$, derivative $LB\xi_t$, or operator-domain
preservation is required.

The actual interaction-picture generator $\lambda(t)B_H(t)$ is bounded
and strongly measurable with integrable norm. Its Dyson propagator exists
on the full Hilbert space. Bounded-generator Duhamel, applied to the
ideal reference vectors only, gives
\begin{equation}
 \|(U_I(\tau)-V(\tau))\psi\|
   \leq2\int_0^\tau |\lambda(t)|\sqrt{2tC}\,dt
   \leq2A_{\rm abs}\sqrt{2\tau C}.
\end{equation}
The actual Schr\"odinger propagator is
$U(\tau)=e^{-i\tau L}U_I(\tau)$. Adding its final free factor, and
using the form bound on $V(\tau)\psi$, yields the finite-pulse estimate
\begin{equation}
 \boxed{\begin{gathered}
 \|(U(\tau)-V(\tau))\psi\|
       \leq(1+2A_{\rm abs})\sqrt{2\tau C}.
 \end{gathered}}
 \label{eq:chvisd-finite-pulse-error}
\end{equation}
This is a state-specific finite-energy estimate. For rectangular pulses,
$L+\lambda B$ is self-adjoint on $D(L)$ by bounded perturbation and has
the same form domain after a finite scalar shift. Its unitary preserves
that form domain. This also makes the endpoint mean energies and bounded
switching payments well defined for the stated input. The proof of
Eq.~\eqref{eq:chvisd-finite-pulse-error} itself uses bounded Dyson evolution
and the reference form estimates, not an assumed invariance of $D(L)$.

\subsection{Exact ideal rotation and phase-free pointer gate}
\label{sec:chvisd-pointer-gate}

Take the ideal source rotation
\begin{equation}
 R=e^{i\pi S/4}.
 \label{eq:chvisd-source-quarter-rotation}
\end{equation}
It is generated by $B_{\rm rot}=S$ with signed area
$\int\lambda_{\rm rot}=-\pi/4$, hence
$A_{\rm rot}=\pi/4$. The full source operator is the original
six-lane Swap with its four unforced directions; it is not replaced by
a bare Hamiltonian during the pulse.

Let $P_{ba}$ be the declared source-only projector onto the ordered
state selected by the receiver, and let $X_{\rm ptr}$ be the pointer
Pauli matrix. Assign the bounded projector generator
\begin{equation}
 B_{\rm gate}=P_{ba}\otimes\frac{I-X_{\rm ptr}}{2}.
 \label{eq:chvisd-projector-pointer-generator}
\end{equation}
At area $\int\lambda_{\rm gate}=\pi$,
\begin{equation}
 \begin{aligned}
 G=e^{-i\pi B_{\rm gate}}
   &=I-2B_{\rm gate}\\
   &=(I-P_{ba})\otimes I+P_{ba}\otimes X_{\rm ptr}.
 \end{aligned}
 \label{eq:chvisd-phase-free-controlled-pointer}
\end{equation}
Thus a blank pointer $|0\rangle$ is flipped precisely on the selected
source subspace, with no extra conditional phase. Both generators have
norm at most one and commute with $H_f$. Their finite source--pointer
matrices act on all $16\times2$ coordinates while the entire field
remains in the dynamics.

For reference, put $P^{\rm ord}_{ab}=|ab\rangle\langle ab|$.
In the ordered prepared-lane basis $|ab\rangle,|ba\rangle$
define $Z=P^{\rm ord}_{ab}-P_{ba}$, $X=S$, and
$Y=-i|ab\rangle\langle ba|+i|ba\rangle\langle ab|$.
Then
\begin{equation}
\begin{gathered}
P_{ba}=\frac{I-Z}{2},\\
 R^\dagger P_{ba}R=\frac{I+Y}{2}.
\end{gathered}
\label{eq:chvisd-pointer-observable-signs}
\end{equation}
The two corresponding ideal signed pointer means are
$\langle Z\rangle=1-2P(1)$ without $R$ and
$\langle Y\rangle=2P(1)-1$ after $R$. These signs belong to this
explicit ordered basis. Any source-current convention used elsewhere
must be translated by its actual source operator map.

\subsection{Two-pulse error, finite switching work and event time}
\label{sec:chvisd-readout-two-pulse-budget}

Let $U_{\rm rot},U_{\rm gate}$ be the actual rectangular pulses and
$R,G$ the ideal ones. Unitary invariance and one hybrid intermediate state
give
\begin{equation}
 \begin{aligned}
 &\|U_{\rm gate}U_{\rm rot}\psi-GR\psi\|\\
 &\quad\leq\|(U_{\rm rot}-R)\psi\|
                +\|(U_{\rm gate}-G)R\psi\|\\
 &\quad\leq r,\\
 r={}&\left(1+\frac\pi2\right)\sqrt{2C\tau_{\rm rot}}
       +(1+2\pi)\sqrt{2C\tau_{\rm gate}}.
 \end{aligned}
 \label{eq:chvisd-two-pulse-readout-error}
\end{equation}
The second pulse is tested on the ideal state $R\psi$, whose $H_f$
expectation is still bounded by $F$. The estimate does not require a new
field-energy estimate on the actual intermediate state. A pointer-only
readout omitting the rotation uses the second term alone.

For any final pointer event, the probability discrepancy is at most $r$,
because pure-state trace distance is no larger than the vector difference.
Partial trace and purification give the same bound for a mixed input.
A binary pointer value in $\{-1,+1\}$ consequently has mean error at
most $2r$. Pointer initialization and detector errors, when present,
must be added using their declared all-outcome probability or joint-state
bounds. These inequalities introduce no additional measurement rule.

If the actual preparation is known only to satisfy
$T(\rho_{\rm act},\rho_{\rm ref})\leq\delta_{\rm prep}$, apply the
same full actual pulse and pointer channel to both states. Trace-distance
contractivity adds at most $\delta_{\rm prep}$ to each output probability
error. The reference pulse estimate therefore gives
\begin{equation}
 |P_{\rm act}(E)-P_{\rm ideal,ref}(E)|
       \leq\delta_{\rm prep}+r.
 \label{eq:chvisd-readout-preparation-contraction}
\end{equation}
No field-energy or form-domain assumption is thereby made on
$\rho_{\rm act}$, which may have a component outside the reference lane.
The $C$ bound and the pulse-drift proof concern the ideal reference
states. Finite source-energy claims for an actual preparation, when
needed for a resource account, require its own supplied energy bound.

\paragraph{Calibrated areas use their actual absolute areas.}
The nominal analyser angle is $\chi=-\pi/2$, corresponding to the source
area $\chi/2=-\pi/4$. If a rectangular rotation has
$|\delta\chi|\leq b$, its frozen-unitary error relative to $R$ is at
most $b/2$, and its absolute area is at most $\pi/4+b/2$. If the pointer
area is $\pi+\delta_g$ with $|\delta_g|\leq u$, its frozen-unitary
error relative to $G$ is at most $u$ and its absolute area is at most
$\pi+u$. A sufficient full bound is therefore
\begin{equation}
 \begin{aligned}
 r_{\rm cal}\leq{}&\frac b2+u\\
   &+\left(1+\frac\pi2+b\right)\sqrt{2C\tau_{\rm rot}}\\
   &+(1+2\pi+2u)\sqrt{2C\tau_{\rm gate}}.
 \end{aligned}
 \label{eq:chvisd-calibrated-pulse-error}
\end{equation}
This follows from $\|e^{-i\delta B}-I\|\leq|\delta|$ for a
self-adjoint contraction and the same hybrid comparison. A net-area
uncertainty does not bound the absolute area of a pulse with uncontrolled
sign reversals; that case must supply $A_{\rm abs}$ itself. Nominal-area
feasibility values use $b=u=0$ in the pulse estimate. A separate timing
or accumulated-phase allowance, for example $|\Omega\,\delta\eta|$,
is not an analyser-axis error unless its declared control map makes it
one. General calibrated work bounds likewise replace the nominal areas
in Eq.~\eqref{eq:chvisd-finite-readout-switch-work} by these actual
absolute-area ceilings.

There are two equivalent ways to allocate the rotation-axis error, and
they must not be added together. Equation~\eqref{eq:chvisd-calibrated-pulse-error}
compares every pulse with the \emph{nominal} ideal rotation; the separate
decoder phase/age error then excludes the same $b$. Alternatively compare
the finite rotation with its \emph{actual frozen} area and include $b$ in
the decoder's ideal-axis phase allowance. The finite-pulse part in that
allocation is
\begin{equation}
 \begin{aligned}
 r_{\rm actual\ axis}\leq{}&u
   +\left(1+\frac\pi2+b\right)\sqrt{2C\tau_{\rm rot}}\\
   &+(1+2\pi+2u)\sqrt{2C\tau_{\rm gate}},
 \end{aligned}
 \label{eq:chvisd-actual-axis-drift-budget}
\end{equation}
with no additional $b/2$. Both allocations use the actual absolute areas
in their drift bounds. A nominal-axis benchmark with $b=0$, whose stated
phase allowance comes entirely from age/timing, has no such axis payment.

For a desired positive vector-error allowance
$r_{\rm rot}+r_{\rm gate}=r_*$, one sufficient finite choice is
\begin{equation}
 \begin{aligned}
 \tau_{\rm rot}
   &\leq\frac{r_{\rm rot}^2}
       {2C^{\rm up}(1+\pi/2)^2},\\
 \tau_{\rm gate}
   &\leq\frac{r_{\rm gate}^2}
       {2C^{\rm up}(1+2\pi)^2}.
 \end{aligned}
 \label{eq:chvisd-sufficient-finite-pulse-times}
\end{equation}
Choose strictly positive times. Their finite rectangular control heights
are $\pi/(4\tau_{\rm rot})$ and $\pi/\tau_{\rm gate}$; tightening the
error allowance charges these larger amplitudes explicitly.

At each on/off jump the changed part of the Hamiltonian is a bounded
operator. The sum of absolute mean switching payments is therefore at
most
\begin{equation}
 \begin{aligned}
 W_{\rm abs}^{\rm switch}
  &\leq2\left(\frac{A_{\rm rot}}{\tau_{\rm rot}}
                     +\frac{A_{\rm gate}}{\tau_{\rm gate}}\right)\\
  &=\frac{\pi}{2\tau_{\rm rot}}
                     +\frac{2\pi}{\tau_{\rm gate}}.
 \end{aligned}
 \label{eq:chvisd-finite-readout-switch-work}
\end{equation}
Directly changing from one rectangular pulse to the next satisfies this
same upper bound by the triangle inequality. This is an absolute bound
on the declared driven switching work, not a deduction of autonomous
controller capital. Initial source/field energy, the supplied blank
pointer, actuator resources, retained outcomes, and any eventual reset
remain separate recipients. Time and energy use the original model
chart; no new SI assignment is supplied.

The selected projector control need not commute with the original
source--field charge. Its actuator and phase/action reference are part
of this external receiving model. The switching-work estimate does not
by itself supply an autonomous charge-conserving actuator or its
microscopic SAM realization.

The state to which the ideal readout refers is the state at pulse entry.
Its ideal $H_f$ bound is used throughout the comparison because the
reference controls preserve $H_f$. The actual classical pointer result
is available only after the latency
$\tau_{\rm rot}+\tau_{\rm gate}$ and its subsequent registered readout.
Calling it the source value at completion would require the corresponding
event-time and phase-advance map. If the complete physical readout must
lie inside the original certified window, choose its entry and positive
pulse durations accordingly. No rate certificate is extended merely by
omitting this finite latency.

The same fixed rotation, pointer gate, durations and output convention are
used for all declared input-age hypotheses $\rho(\eta_j)$. The unknown
label is not supplied to a control schedule. Its physical preparation and
the two-arm arrival map are separate receiving-model data; a
label-dependent ideal timing choice cannot be substituted for a single
implemented decoder.

